\addcontentsline{toc}{subsubsection}{Solution 28. The Degree-Normalized Cubic Squeeze}

\begin{dlsummary}[title={\bfseries Solution 28. The Degree-Normalized Cubic Squeeze},breakable]

\begin{proof}
Define $g(x)=f(x)/x^3$, so $f(x)=x^3g(x)$. Substitution into the
given inequalities gives
\[
xg(x)+yg(y)\leq(x+y)g(x+y),
\qquad
g(x)+g(y)\geq2g(x+y).
\]
Equivalently,
\[
\frac{xg(x)+yg(y)}{x+y}
\leq g(x+y)\leq\frac{g(x)+g(y)}2
\qquad(x,y>0). \tag{1}
\]

\medskip
\noindent
\textbf{Step 1. Establish dyadic invariance and monotonicity.}

Taking $y=x$ in (1) yields
\[
g(x)\leq g(2x)\leq g(x),
\]
so
\[
g(2x)=g(x)\qquad(x>0). \tag{2}
\]
The lower and upper bounds in (1) also imply
\[
2\bigl(xg(x)+yg(y)\bigr)\leq(x+y)\bigl(g(x)+g(y)\bigr),
\]
which simplifies to
\[
(x-y)\bigl(g(x)-g(y)\bigr)\leq0.
\]
Thus $g$ is non-increasing.

\medskip
\noindent
\textbf{Step 2. Show that $g$ is constant.}

Fix $a>0$. For any $x\in[a,2a]$, monotonicity and (2) give
\[
g(a)=g(2a)\leq g(x)\leq g(a),
\]
so $g(x)=g(a)$ throughout $[a,2a]$. These intervals cover
$(0,\infty)$ as $a$ varies, and therefore $g$ is constant: there is
some $c\in\mathbb R$ such that $g(x)=c$ for all $x>0$.
Consequently,
\[
f(x)=cx^3\qquad(x>0).
\]

\medskip
\noindent
\textbf{Verification.}

For any real constant $c$, substituting $f(x)=cx^3$ into the first
inequality reduces it to
$c(x+y)\leq c(x+y)$, and the second to
$2c\geq2c$. Hence every such function satisfies both conditions.
\end{proof}

\end{dlsummary}